% !PRINTED 528
\weberpage{528}

Eliminiren wir aber aus den Gleichungen (2) und (3) die Potenzen von $\Theta$, so ergiebt sich eine Relation von der Form
\weq{6}{\omega = k_1\omega_1 + k_2\omega_2 + \cdots + k_n\omega_n.}

Man kann also jede Zahl und jedes Functional des Körpers $\Omega$ in der Form (6) darstellen, wenn man für $k_1, k_2, \ldots, k_n$ rationale Zahlen oder Functionale setzt. Diese Darstellung ist für ein gegebenes $\omega$ [wegen (5)] nur auf eine Art möglich, und $\omega$ ist eine Zahl, wenn die $k_1, k_2, \ldots, k_n$ Zahlen sind, dagegen ein Functional, wenn unter den $k$ auch Functionale vorkommen.

Ein solches System von Zahlen, wie
\[
\omega_1, \omega_2, \ldots, \omega_n,
\]
nennen wir eine \gesperrt{Basis des Körpers $\Omega$}. Eine solche Basis bilden auch die Potenzen von $\Theta$:
\[
1, \Theta, \Theta^2, \ldots, \Theta^{n-1}.
\]

Bezeichnen wir mit $\omega_{r,1}, \omega_{r,2}, \ldots, \omega_{r,n}$ die mit $\omega_r$ conjugirten Zahlen, so ist das Determinantenquadrat
\weq{7}{\Delta(\omega_1, \omega_2, \ldots, \omega_n) = (\Sigma \pm \omega_{1,1}\omega_{2,2}\ldots\omega_{n,n})^2}
eine symmetrische Function der conjugirten Werthe $\Theta_1, \Theta_2, \ldots, \Theta_n$ und ist folglich eine rationale Zahl.

Diese Zahl heisst die \gesperrt{Discriminante des Systemes} $\omega_1, \omega_2, \ldots, \omega_n$. Insbesondere ist
\weq{8}{\Delta(1, \Theta, \Theta^2, \ldots, \Theta^{n-1}) = \begin{vmatrix} 1, & \Theta_1, & \ldots, & \Theta_1^{\,n-1}\\ 1, & \Theta_2, & \ldots, & \Theta_2^{\,n-1}\\ \multicolumn{4}{c}{\cdots\cdots\cdots\cdots\cdots\cdots}\\ 1, & \Theta_n, & \ldots, & \Theta_n^{\,n-1} \end{vmatrix}^2}
die Discriminante der Gleichung (1) (Bd.\ I, \S.\,46), für die man auch, wenn $N$ das Zeichen für die Norm ist,
\weq{9}{(-1)^{\frac{n(n-1)}{2}} N f'(\Theta)}
setzen kann. Diese Zahl ist also sicher von Null verschieden. Nach dem Multiplicationssatze der Determinanten (Bd.\ I, \S.\,27) ergiebt sich aus (3) und (7):
\weq{10}{\Delta(\omega_1, \omega_2, \ldots, \omega_n) = H^2\Delta(1, \Theta, \Theta^2, \ldots, \Theta^{n-1}),}
woraus man schliesst, dass die Discriminante einer Basis von $\Omega$ immer von Null verschieden ist. Da $H$ eine rationale Zahl ist, so folgt, dass das Verhältniss der Discriminanten verschiedener Basen das Quadrat einer rationalen Zahl ist, und dass die
%CONT>>
